% Vietnamese translation. Original: Carreira, Amaral, Vangheluwe (eds.), 2020.
% Source: PDF page 5 (dedication).
\chapter*{Lời tưởng niệm}
Tưởng nhớ người bạn và đồng nghiệp thân thiết Bernhard Schätz (1965--2017), một nhà khoa học xuất sắc và một trong những người sáng lập mạng lưới MPM4CPS. Chúng tôi sẽ mãi nhớ đến tinh thần làm việc đáng quý của ông, nguồn cảm hứng cho mọi người xung quanh. Chúng tôi vô cùng nhớ sự nhiệt tình và thái độ luôn tích cực của ông.

% Source: printed vii / PDF 6.
\chapter*{Lời giới thiệu}
\addcontentsline{toc}{chapter}{Lời giới thiệu}
Sự gia tăng mạnh mẽ về năng lực tính toán của máy tính để bàn, bắt đầu từ những năm 1980, đã đưa mô hình hóa bằng tính toán đến với số lượng kỹ sư và nhà khoa học tăng theo cấp số nhân. Tuy nhiên, một trở ngại lớn đối với việc khai thác tiềm năng này là yêu cầu phải biểu diễn mô hình bằng các ngôn ngữ lập trình, chẳng hạn FORTRAN và Algol. Trong khi các tổ chức như phòng thí nghiệm quốc gia có thể thuê kỹ sư phần mềm để lập trình các mô hình toán học, người dùng máy tính để bàn lại không có những nguồn lực tương tự.

Theo tiến trình tự nhiên, những ngôn ngữ mô hình hóa bằng tính toán dễ tiếp cận hơn đã được phát triển, phù hợp với thực tiễn mô hình hóa và các hệ hình thức trong nhiều lĩnh vực. Chẳng hạn, các ngôn ngữ thuật toán dựa trên mảng như MATLAB\textsuperscript{\textregistered} tỏ ra đặc biệt hữu ích trong xử lý tín hiệu. Trong điều khiển, mô hình thường được biểu diễn dưới dạng sơ đồ khối. Trong truyền thông, mạng Petri là một cách biểu diễn phù hợp cho các giao thức. Các hệ vật lý có thể được biểu diễn thuận tiện nhất bằng mô hình phi nhân quả, chẳng hạn đồ thị liên kết năng lượng (bond graph). Bên cạnh đó, các môi trường mô hình hóa trực quan, có vai trò không kém phần quan trọng, cũng được phát triển, thường dựa trên công nghệ giao diện người dùng đồ họa (GUI) đang hình thành.

Việc phổ cập mô hình hóa bằng tính toán nhờ máy tính để bàn ban đầu đem lại công cụ mô phỏng hỗ trợ công việc của kỹ sư và nhà khoa học. Theo thời gian, các sản phẩm tính toán này có thêm nhiều giá trị. Mô hình hóa và mô phỏng bằng tính toán không chỉ trở thành bước quan trọng trong hoạt động phân tích của từng cá nhân; việc chia sẻ mô hình giữa các nhóm, phòng ban và thậm chí giữa các tổ chức cũng ngày càng tạo ra giá trị lớn hơn.

Vào khoảng bước chuyển sang thiên niên kỷ mới, xu hướng xem mô hình là sản phẩm có vai trò trung tâm trong thiết kế các hệ thống kỹ thuật bắt đầu định hình. Khi còn làm việc tại DLR Oberpfaffenhofen, tôi và Hans Vangheluwe gặp nhau tại một cuộc họp thiết kế ngôn ngữ Modelica. Những cuộc trao đổi của chúng tôi dần làm rõ các vấn đề nghiên cứu nền tảng cần giải quyết để đáp ứng nhu cầu ngày càng tăng về chia sẻ và tái sử dụng mô hình. Thách thức lớn đặt ra là kết nối, kết hợp và tích hợp những mô hình khác nhau về cú pháp (chẳng hạn dạng văn bản, các yếu tố đồ họa khác nhau và các kiểu quan hệ khác nhau), ngữ nghĩa (mệnh lệnh hay khai báo, luồng dữ liệu hay luồng điều khiển, có yếu tố thời gian hay không xét thời gian), bối cảnh sử dụng (cá nhân hay nhóm, sản phẩm thiết kế then chốt hay một thử nghiệm riêng lẻ, gỡ lỗi hay vận hành thực tế), góc nhìn (chức năng hay hiệu năng, cấu trúc hay hành vi), mục đích (phân tích hay tổng hợp, an toàn hay sáng tạo), quy trình (từ toán học đến tính toán hay từ toán học đến ngữ nghĩa tính toán, mô hình hóa từ các nguyên lý cơ bản hay dựa trên dữ liệu), và luồng công việc (từ trên xuống hay từ dưới lên, tích hợp liên tục hay chia sẻ định kỳ, tinh chỉnh liên tục hay kiểm thử đối chiếu).

Một ngày nọ, tôi cùng Georg Gruebel, người phụ trách nhóm của tôi tại DLR Oberpfaffenhofen, bàn về việc tổ chức một phiên chuyên đề tại Hội nghị chuyên đề quốc tế IEEE về thiết kế hệ thống điều khiển có máy tính hỗ trợ (CACSD) năm 2000. Chúng tôi chưa tìm được tên phù hợp cho phiên này, và ông ra về. Khoảng mười lăm phút sau, ông trở lại văn phòng của tôi với một cái tên cho lĩnh vực đang hình thành: Mô hình hóa đa mô thức được tự động hóa bằng máy tính (Computer Automated Multiparadigm Modeling, CAMPaM).

Sáng kiến tại CACSD 2000 thành công đến mức phải tổ chức hai phiên hội nghị. Sau đó, chúng tôi tiếp tục thực hiện các số chuyên đề trên SCS SIMULATION và IEEE Transactions on Control System Technology. Năm 2004, chúng tôi khởi động chuỗi hội thảo Bellairs về CAMPaM, vẫn đang hoạt động đến nay.

Theo thời gian, mục tiêu tự động hóa bằng máy tính tỏ ra quá lạc quan. Chúng tôi điều chỉnh kỳ vọng bằng cách mở rộng phạm vi sang mô hình hóa đa mô thức (MPM) nói chung, bao gồm cả các cách tiếp cận có máy tính hỗ trợ. Theo đó, một chuỗi hội thảo về MPM đã được tổ chức trong khuôn khổ Hội nghị quốc tế ACM/IEEE về ngôn ngữ và hệ thống kỹ nghệ hướng mô hình (MODELS) từ năm 2010 đến năm 2015.

% Source: printed viii / PDF 7.
Ngày nay, nhận định có tính tiên liệu của chúng tôi vào thời điểm chuyển giao thiên niên kỷ được thể hiện qua nhu cầu cấp thiết về mô hình khi các hệ thống thực-ảo (Cyber-Physical Systems, CPS) ngày càng phổ biến. CPS đặt tính toán và truyền thông vào vị trí trung tâm của hoạt động vật lý. Đặc biệt, điều này tạo ra mức độ linh hoạt hay tính mở khác biệt rõ rệt so với cách vận hành tính toán cứng nhắc, đơn nhất trong mô thức hệ thống nhúng có trước CPS. CPS mở theo chiều dọc: nền tảng tính toán, cả phần cứng lẫn phần mềm, có thể thay đổi độc lập với logic mà các phép tính thực hiện. CPS cũng mở theo chiều ngang: các hệ thống kết nối với những hệ thống khác để đạt mục tiêu tổng thể thông qua hành vi xuất hiện sau khi triển khai. Chính xu hướng ngày càng mở này đòi hỏi sử dụng mô hình trên nhiều mô thức ngày một nhiều hơn, cả trong thiết kế hệ thống lẫn trong quá trình vận hành, tức tại thời điểm chạy.

Vì vậy, việc biên soạn giáo trình \emph{Cơ sở mô hình hóa đa mô thức cho các hệ thống thực-ảo} là rất đúng lúc. Nhiều nỗ lực nghiên cứu đã được củng cố qua các năm, trong khi nhu cầu thực tiễn về mô hình hóa đa mô thức vẫn tiếp tục tăng. Cuốn sách tập hợp những nội dung có giá trị về các khía cạnh chủ yếu của CPS, từ mô hình hóa các hệ vật lý, tính toán và truyền thông đến điều khiển, các đặc tính hiệu năng và kiến trúc hệ thống. Là một giáo trình nền tảng về mô hình hóa đa mô thức cho CPS, cuốn sách đặt tất cả những khía cạnh này trên cơ sở của mô hình hóa đa mô thức.

Trong nhiều năm qua, tôi rất vui được làm việc với tất cả các tác giả và biên tập viên của tuyển tập này. Cuốn sách cho thấy lĩnh vực CAMPaM đang trưởng thành, và điều đó khiến tôi rất hào hứng. Xin cảm ơn tất cả những người đã dành thời gian và công sức để làm nên cuốn sách!

\begin{flushright}
Natick, Massachusetts, Hoa Kỳ, tháng 5 năm 2019\\
\textbf{Pieter J. Mosterman}\\
Nhà khoa học nghiên cứu trưởng và Giám đốc\\
Văn phòng Nghiên cứu và Công nghệ tiên tiến của MathWorks
\end{flushright}

% Source: printed ix / PDF 8.
\chapter*{Lời nói đầu}
\addcontentsline{toc}{chapter}{Lời nói đầu}
Mô hình hóa và phân tích là những hoạt động thiết yếu trong phát triển hệ thống thực-ảo. Tính liên ngành vốn có của CPS đòi hỏi sử dụng những kỹ thuật mô hình hóa khác nhau thuộc các ngành khác nhau. Đồng thời, để các chuyên ngành có thể trao đổi với nhau, cần có kiến thức nền tảng chung về các hệ hình thức.

Một nhận thức quan trọng là các mô hình riêng biệt phải được liên kết nhất quán để tạo thành biểu diễn đầy đủ của hệ thống. Nhờ đó, có thể thực hiện những hoạt động ở cấp độ toàn hệ thống như phân tích hiệu năng, mô phỏng vét cạn và kiểm chứng, mô phỏng có phần cứng trong vòng lặp, xác định bộ tham số tổng thể tốt nhất cho hệ thống, xây dựng nguyên mẫu và hiện thực hóa. Có thể hình dung một hệ hình thức thống nhất biểu diễn các mối liên hệ giữa những mô hình riêng biệt nêu trên, qua đó cho phép suy luận về chúng. Tương tự, những tiến bộ gần đây cho thấy việc kết nối các ngành bằng đồng mô phỏng giúp chúng phối hợp mà không buộc phải sử dụng công cụ hay phương pháp thiết kế mới.

Trước hết, CPS là một lĩnh vực tương đối mới, chưa có các tiêu chuẩn đã được xác lập, với nhiều bên tham gia; mỗi bên lại có định nghĩa và cách tiếp cận hơi khác nhau. Mô hình hóa đa mô thức bao quát một phổ rộng các chủ đề khác nhau và có định nghĩa khái quát, linh hoạt, chưa mang tính hình thức. Vì vậy, các công cụ hiện có được xây dựng theo nhiều hướng khác nhau, trong một lĩnh vực trải dài từ các cách tiếp cận rất thiên về lý thuyết đến các cách tiếp cận hướng ứng dụng.

Mục tiêu của cuốn sách là cung cấp một tài liệu tham khảo tập hợp nhất quán những kiến thức đã được xác lập về cơ sở và các hệ hình thức dùng để mô hình hóa CPS. Cuốn sách tập trung vào những kiến thức và thực hành vững chắc, tiên tiến, hay các nội dung nền tảng, thay vì trình bày nghiên cứu; đồng thời cố gắng xử lý thống nhất tính liên ngành của mô hình hóa CPS. Sách có thể làm cầu nối cho người bước vào lĩnh vực CPS từ những ngành kỹ thuật có liên quan, chẳng hạn kỹ thuật điện, cơ học, vật lý hay khoa học máy tính. Dù chủ yếu dành cho học viên cao học và sau đại học thuộc khoa học máy tính và kỹ thuật, sách cũng có thể được sử dụng rộng rãi làm tài liệu tham khảo cho người làm thực tiễn.

% Source: printed xi / PDF 9.
\chapter*{Lời cảm ơn}
\addcontentsline{toc}{chapter}{Lời cảm ơn}
Cuốn sách nhận được sự hỗ trợ của COST Action IC1404, Mô hình hóa đa mô thức cho các hệ thống thực-ảo (MPM4CPS). COST được Chương trình khung Horizon 2020 của Liên minh châu Âu hỗ trợ.
